\chapter{Otoacoustic Emissions (OAE) Test}

\section*{What Is This Test?}
An otoacoustic-emissions test measures faint sounds produced by healthy outer hair cells in the cochlea in response to an inserted sound stimulus. It is a quick, objective hearing test used widely for newborn hearing screening and for selected follow-up hearing evaluations.

\section*{Alternative Names}
OAE Test; Otoacoustic Emission Test; Newborn OAE Hearing Screen; Evoked Otoacoustic Emissions; TEOAE; DPOAE

\section*{Why It Is Ordered}
Universal newborn hearing screening, follow-up after a failed or incomplete screen, evaluation of suspected cochlear hearing loss when behavioral responses are unreliable, and monitoring of selected ototoxic exposures or ear-related conditions.

\section*{How This Test Is Performed}
A small probe is placed gently in the ear canal. The probe presents tones or clicks and measures the echo generated by the cochlea. The test is painless and usually takes only a few minutes per ear; a quiet, still environment or natural sleep improves the recording.

\section*{Preparation, Risks, and Limitations}
Excess earwax, middle-ear fluid, background noise, movement, or a poor probe seal can produce a refer result. The test is noninvasive and has no meaningful physical risk. A pass or refer result is a screening outcome, not a complete diagnosis; infants who do not pass require timely diagnostic audiology, which may include ABR.

\section*{References}
\begin{enumerate}
  \item MedlinePlus. Hearing Tests for Children and Adults.
  \item Joint Committee on Infant Hearing. Year 2019 Position Statement.
\end{enumerate}

\clearpage
\section*{Understanding Results and Follow-up}
The laboratory report should identify the specimen, method, units, reference interval or decision limit, and whether the result is positive, negative, abnormal, or inconclusive. Reference intervals vary with age, sex, pregnancy, specimen type, assay, and laboratory; a result outside a range is not automatically a diagnosis, and a result within a range does not always exclude disease. Discuss unexpected results with the ordering clinician, who can decide whether repeat or confirmatory testing, treatment, or referral is appropriate. Do not change medicines or delay urgent care based on an isolated result.


\section*{Regulatory Status and Authorization Date}
This chapter describes a test category, not a specific commercial product. FDA, CDSCO, EU IVDR, and MHRA approval or registration dates therefore cannot be assigned without the assay manufacturer, product name, instrument, and intended use. For laboratory-developed tests, an individual agency approval date may not exist. See the Preface for the interpretation of regulatory status.

\clearpage
